% Risk Parity from Scratch — research-note PDF.
%
% PROSE, CODE AND TABLES ARE VERBATIM from the website article:
%   site/src/content/articles/risk-parity-futures.tsx
%   site/src/content/articles/rp-code.ts
% Metadata from lib/articles.ts; project-card fallbacks from lib/topics.ts.
%
% CHART generated by figures/make_risk_parity_futures.py from the article's own
% data module (data/risk-parity-from-scratch.ts) — the numbers the site's <Bar>
% renders. See that script for a note on the article's own inconsistency: its
% chart data covers four ETFs while its tables discuss five futures.
\documentclass{pyportfolios-note}

\noteshorttitle{Risk Parity from Scratch}

\begin{document}
\thispagestyle{notecover}

\notemasthead{Quantitative Insights}{Tutorial}{Q3 2026}{}

\notetitle{Risk Parity\\from Scratch}

\notedeck{Balancing a portfolio by risk rather than by capital — and why futures
are the natural vehicle.}

\notelede{A traditional 60/40 portfolio looks balanced by \textit{capital} — but
not by \textit{risk}. Because equities are \textasciitilde3× as volatile as
bonds, that 60/40 actually derives \textasciitilde90\% of its risk from stocks.
\textbf{Risk parity} flips the question: instead of asking how much money each
asset gets, it asks how much \textit{risk} each contributes — and equalizes it.
Futures are the natural vehicle, since risk parity leans on leverage to lift
low-vol assets to their weight.}

\notecard
  {Portfolio Optimization}
  {NumPy · Pandas · SciPy · Riskfolio-lib}
  {ES, ZN, GC, HG, CL futures}
  {Jan 2015 – Dec 2024}
  {Balancing a multi-asset book by risk contribution rather than capital}

% ============================================================= 1 Summary =====
\noteneedroom{150bp}
\notesection{1.\notenumsep Summary}

\begin{multicols}{2}
Every asset in a portfolio contributes some share of the total risk — and those
shares are rarely equal. Risk parity finds the weights that make each asset's
\textbf{risk contribution} identical, so no single position dominates the
portfolio's fate. Because low-volatility assets (bonds) need large weights to
pull their risk-weight up, risk parity portfolios are typically \textbf{levered}
to reach an equity-like return — which is exactly why they're built with
futures.
\end{multicols}

% =========================================================== 2 Intuition =====
\noteneedroom{170bp}
\notesection{2.\notenumsep Intuition}

\begin{multicols}{2}
\notesubsection{2.1\notenumsep Capital weight $\neq$ risk weight}
Put 60\% in stocks (15\% vol) and 40\% in bonds (5\% vol). The stock sleeve's
risk is 60\%×15\% = 9 ``units''; the bond sleeve's is 40\%×5\% = 2. Stocks
supply \textasciitilde80–90\% of the total. Your ``balanced'' portfolio is a
stock portfolio wearing a bond costume.

\notesubsection{2.2\notenumsep Equalize the risk, not the dollars}
Risk parity down-weights the volatile assets and up-weights the calm ones until
each contributes the same risk. The equity slice shrinks; bonds, gold and rates
grow. The portfolio finally becomes diversified in the way that matters — by
\textit{source of risk}.

\notesubsection{2.3\notenumsep Leverage is a feature, not a bug}
A risk-balanced portfolio is dominated by low-vol assets, so its raw volatility
(and return) is low. To reach an equity-like target you scale the whole thing up
with leverage — cheap and clean in futures. This is the mechanism behind ``All
Weather''-style funds.
\end{multicols}

% ================================================= 3 Theory & Mechanics ======
\noteneedroom{180bp}
\notesection{3.\notenumsep Theory \& Mechanics}

For weights $w$ and covariance $\Sigma$, portfolio volatility is
$\sigma_p = \sqrt{w^\top\Sigma w}$. Each asset's \textbf{marginal risk
contribution} is:

\[ MRC_i = \frac{\partial \sigma_p}{\partial w_i} = \frac{(\Sigma w)_i}{\sigma_p},
   \qquad RC_i = w_i \cdot MRC_i \]

and the $RC_i$ sum exactly to $\sigma_p$ (Euler's theorem). \textbf{Risk parity}
seeks weights where all $RC_i$ are equal:

\[ RC_i = \frac{\sigma_p}{n} \quad \forall i \]

There's no closed form, so we minimize the dispersion of risk contributions
numerically — first by hand in SciPy, then cross-checked with
\textbf{Riskfolio-lib}.

% ==================================== 4 Applied Example — Five Futures =======
\noteneedroom{200bp}
\notesection{4.\notenumsep Applied Example — Five Futures}

\notesubsection{4.1\notenumsep A cross-asset futures universe}

Five liquid futures spanning the major risk factors — equities, rates, and three
commodities:

\begin{notetable}{@{}p{100bp}p{200bp}p{196bp}@{}}
\thcell{Ticker} & \thcell{Contract} & \thcell{Class} \\
\theadrule
ES=F & S\&P 500 E-mini & Equity \\ \rowrule
ZN=F & 10y T-Note & Rates \\ \rowrule
GC=F & Gold & Metal (haven) \\ \rowrule
HG=F & Copper & Metal (cyclical) \\ \rowrule
CL=F & WTI Crude & Energy \\
\end{notetable}

\begin{notecode}
TICKERS = ["ES=F", "ZN=F", "GC=F", "HG=F", "CL=F"]
NAMES = {"ES=F": "S&P 500", "ZN=F": "10y Note", "GC=F": "Gold", "HG=F": "Copper", "CL=F": "WTI"}
START, END = "2015-01-01", "2024-12-31"

def load_prices(tickers, start, end):
    """Continuous futures adjusted-close via yfinance."""
    import yfinance as yf
    df = yf.download(tickers, start=start, end=end, auto_adjust=True, progress=False)["Close"]
    return df[tickers].dropna()

px = load_prices(TICKERS, START, END)
rets = px.pct_change().dropna()

Sigma = rets.cov() * 252
vol = pd.Series(np.sqrt(np.diag(Sigma)), index=TICKERS)
print(f"{len(px)} trading days, {px.index[0].date()} → {px.index[-1].date()}\n")
print("Annualised volatility:")
print(vol.rename(index=NAMES).round(3))
\end{notecode}

\begin{notetable}{@{}p{240bp}>{\raggedleft\arraybackslash}p{256bp}@{}}
\thcell{Contract} & \thcell{Ann. volatility} \\
\theadrule
S\&P 500 & 17.9\% \\ \rowrule
10y Note & 5.5\% \\ \rowrule
Gold & 14.7\% \\ \rowrule
Copper & 22.1\% \\ \rowrule
WTI & 114.9\% \\
\end{notetable}

\notecallout{Data note}{Continuous \texttt{=F} series include roll gaps — and
WTI's eye-watering 114.9\% vol is dominated by the April 2020 episode, when the
front-month contract briefly traded \textit{negative} and daily percentage
returns exploded. Fine for illustrating risk-parity mechanics (it makes the
concentration problem vivid); production would use a back-adjusted series.}

\notesubsection{4.2\notenumsep The problem with equal weight}

Start naive: put 20\% in each future. Watch how unequal the \textit{risk}
contributions are — crude alone supplies 87\% of the portfolio's risk, while the
10y note contributes essentially nothing:

\begin{notecode}
def risk_contributions(w, Sigma):
    """Risk contribution of each asset (sums to portfolio vol)."""
    w = np.asarray(w)
    sigma_p = np.sqrt(w @ Sigma @ w)
    mrc = (Sigma @ w) / sigma_p
    return w * mrc, sigma_p

n = len(TICKERS)
w_eq = np.repeat(1/n, n)
rc_eq, sig_eq = risk_contributions(w_eq, Sigma.values)

tbl = pd.DataFrame({"weight": w_eq, "risk contrib": rc_eq,
                    "risk share": rc_eq / sig_eq}, index=[NAMES[t] for t in TICKERS])
print(tbl.round(3))
print(f"\nPortfolio vol: {sig_eq:.1%} — but risk shares range "
      f"{(rc_eq/sig_eq).min():.0%} to {(rc_eq/sig_eq).max():.0%}, far from equal.")
\end{notecode}

\begin{notetable}{@{}p{160bp}>{\raggedleft\arraybackslash}p{100bp}>{\raggedleft\arraybackslash}p{110bp}>{\raggedleft\arraybackslash}p{110bp}@{}}
\thcell{Contract} & \thcell{Weight} & \thcell{Risk contrib} & \thcell{Risk share} \\
\theadrule
S\&P 500 & 20\% & 0.011 & 4.4\% \\ \rowrule
10y Note & 20\% & −0.000 & −0.0\% \\ \rowrule
Gold & 20\% & 0.006 & 2.4\% \\ \rowrule
Copper & 20\% & 0.016 & 6.3\% \\ \rowrule
WTI & 20\% & 0.219 & 86.9\% \\
\end{notetable}

Portfolio vol is 25.2\%, and the risk shares range from −0\% to 87\% — ``equal
weight'' is a crude-oil bet in disguise. (The note's share is slightly
\textit{negative}: its co-movement with the rest actually subtracts risk.)

\notesubsection{4.3\notenumsep Solve risk parity from scratch (SciPy)}

Minimize the squared dispersion of risk contributions. At the optimum, every
asset supplies exactly 1/n of the total risk.

\begin{notecode}
from scipy.optimize import minimize

def rp_objective(w, Sigma):
    rc, sigma_p = risk_contributions(w, Sigma)
    target = sigma_p / len(w)
    return np.sum((rc - target) ** 2)          # equalise risk contributions

cons = ({"type": "eq", "fun": lambda w: w.sum() - 1},)
bnds = tuple((0.0, 1.0) for _ in range(n))
res = minimize(rp_objective, w_eq, args=(Sigma.values,), method="SLSQP",
               bounds=bnds, constraints=cons, tol=1e-12)

w_rp = res.x
rc_rp, sig_rp = risk_contributions(w_rp, Sigma.values)

tbl = pd.DataFrame({"RP weight": w_rp, "risk share": rc_rp / sig_rp},
                   index=[NAMES[t] for t in TICKERS])
print(tbl.round(3))
print(f"\nRisk shares now range {(rc_rp/sig_rp).min():.1%} to {(rc_rp/sig_rp).max():.1%} "
      f"— all ≈ {1/n:.0%}. Portfolio vol: {sig_rp:.1%}")
\end{notecode}

\begin{notetable}{@{}p{200bp}>{\raggedleft\arraybackslash}p{140bp}>{\raggedleft\arraybackslash}p{140bp}@{}}
\thcell{Contract} & \thcell{RP weight} & \thcell{Risk share} \\
\theadrule
S\&P 500 & 15.6\% & 20.0\% \\ \rowrule
10y Note & 55.2\% & 20.0\% \\ \rowrule
Gold & 15.7\% & 20.0\% \\ \rowrule
Copper & 11.1\% & 20.0\% \\ \rowrule
WTI & 2.4\% & 20.0\% \\
\end{notetable}

All five risk shares land on exactly 20\%. The allocation inverts the vol
ranking: the calm 10y note gets 55\% of the capital, wild WTI just 2.4\% — and
portfolio vol drops from 25.2\% to 6.9\%.

\notefigure{figures/risk-parity-futures/fig1-risk-share.png}
  {Figure 4.3 · Capital weight vs risk share — equal weight vs risk parity}
  {Equal weighting gives every asset the same capital but very different risk
   shares; risk parity equalises them}
  {Source: Author calculations, from the article's computed results.}

\notesubsection{4.4\notenumsep Validate with Riskfolio-lib}

Our hand-rolled solution should match the library's dedicated risk-parity
optimizer to within solver tolerance — and it does, to 7.5 × 10\textsuperscript{−6}:

\begin{notecode}
import riskfolio as rp

port = rp.Portfolio(returns=rets)
port.assets_stats(method_mu="hist", method_cov="hist")
w_lib = port.rp_optimization(model="Classic", rm="MV", rf=0, b=None)

compare = pd.DataFrame({
    "from scratch (SciPy)": w_rp,
    "Riskfolio-lib": w_lib["weights"].values,
}, index=[NAMES[t] for t in TICKERS])
compare["abs diff"] = (compare.iloc[:, 0] - compare.iloc[:, 1]).abs()
print(compare.round(4))
print(f"\nMax weight difference: {compare['abs diff'].max():.2e}  (solver tolerance)")
\end{notecode}

\begin{notetable}{@{}p{180bp}>{\raggedleft\arraybackslash}p{150bp}>{\raggedleft\arraybackslash}p{150bp}@{}}
\thcell{Contract} & \thcell{From scratch (SciPy)} & \thcell{Riskfolio-lib} \\
\theadrule
S\&P 500 & 0.1557 & 0.1557 \\ \rowrule
10y Note & 0.5523 & 0.5523 \\ \rowrule
Gold & 0.1571 & 0.1571 \\ \rowrule
Copper & 0.1106 & 0.1106 \\ \rowrule
WTI & 0.0243 & 0.0243 \\
\end{notetable}

\notesubsection{4.5\notenumsep Adding leverage to hit a volatility target}

Risk parity's raw vol is low (bonds dominate). Scale the whole book to a target —
say 10\% annualised — with a single leverage multiplier. In futures this costs
only margin, not capital.

\begin{notecode}
target_vol = 0.10
leverage = target_vol / sig_rp
w_levered = w_rp * leverage

print(f"Unlevered RP vol: {sig_rp:.1%}")
print(f"Leverage to reach {target_vol:.0%}: {leverage:.2f}x")
print(f"Gross exposure: {w_levered.sum():.2f}  (vs 1.00 unlevered)\n")
print(pd.Series(w_levered, index=[NAMES[t] for t in TICKERS], name="levered weight").round(3))
\end{notecode}

\begin{notetable}{@{}p{180bp}>{\raggedleft\arraybackslash}p{150bp}>{\raggedleft\arraybackslash}p{150bp}@{}}
\thcell{Contract} & \thcell{Unlevered weight} & \thcell{Levered weight} \\
\theadrule
S\&P 500 & 15.6\% & 22.6\% \\ \rowrule
10y Note & 55.2\% & 80.3\% \\ \rowrule
Gold & 15.7\% & 22.8\% \\ \rowrule
Copper & 11.1\% & 16.1\% \\ \rowrule
WTI & 2.4\% & 3.5\% \\
\end{notetable}

Unlevered vol is 6.9\%, so hitting 10\% takes \textbf{1.45×} leverage — gross
exposure 1.45 instead of 1.00. The note position alone is now 80\% of capital:
this is what ``levered bonds'' means in the risk-parity debate.

% ========================================================== 5 Conclusion =====
\noteneedroom{330bp}
\notesection{5.\notenumsep Conclusion}

\begin{multicols}{2}
\notesubsection{5a.\notenumsep Strengths}
\begin{notebullets}
\item \textbf{True diversification} — balances the \textit{sources} of risk, not
      just the dollars
\item \textbf{No return forecasts needed} — depends only on the covariance
      matrix, sidestepping MVO's biggest weakness
\item \textbf{Robust, stable weights} — small input changes barely move the
      allocation
\item \textbf{Futures-native} — leverage to a vol target is cheap and clean in
      the futures market
\end{notebullets}

\notesubsection{5b.\notenumsep Weaknesses \& Limitations}
\begin{notebullets}
\item \textbf{Leverage brings its own risks} — funding cost, margin calls, and
      forced deleveraging in a crisis (March 2020 hit risk-parity funds hard)
\item \textbf{Covariance is still estimated} — correlations that break down in
      stress undermine the whole premise
\item \textbf{Bond-heavy by construction} — vulnerable when rates and equities
      fall \textit{together} (2022)
\item \textbf{Ignores expected returns} — equalizing risk is agnostic about
      where return actually comes from
\end{notebullets}

\notesubsection{5c.\notenumsep Applications in Practice}
\begin{notebullets}
\item The engine behind ``All Weather'' / risk-parity funds (Bridgewater and
      imitators)
\item A diversification overlay for multi-asset and CTA portfolios
\item A return-forecast-free benchmark to test whether your alpha views actually
      add value
\end{notebullets}

\notesubsection{5d.\notenumsep Alternatives \& Extensions}
\begin{notebullets}
\item \textbf{Hierarchical Risk Parity (HRP)} — clusters assets first, avoiding
      matrix inversion for stabler weights
\item \textbf{Risk budgeting} — the general case, assigning \textit{unequal}
      target risk shares by conviction
\item \textbf{Mean-variance \& Black-Litterman} — the return-driven alternatives
      (here and here)
\item \textbf{CVaR risk parity} — equalize tail-risk contributions instead of
      variance contributions
\end{notebullets}
\end{multicols}

\end{document}
